\subsection{Objetivos}

\par A seguir são apresentados os objetivos gerais e específicos deste Trabalho de Conclusão de Curso.

\subsubsection{Objetivo Geral}

\par Desenvolver um sistema supervisório web com arquitetura orientada a serviços, capaz de apoiar o monitoramento e a supervisão de processos industriais, com serviços bem definidos, reutilizáveis e acessíveis via interfaces web.

\subsubsection{Objetivos Específicos}

\begin{itemize}
    \item Revisar os conceitos de Indústria~4.0, monitoramento industrial, arquitetura orientada a serviços (SOA) e serviços Web, de modo a embasar a solução proposta.
    \item Especificar os requisitos funcionais e não funcionais e a arquitetura lógica do sistema, com destaque para a separação entre camada de apresentação (web), camada de serviços e camada de obtenção ou persistência dos dados de monitoramento.
    \item Implementar os serviços de backend responsáveis pela leitura periódica ou sob demanda, pelo registro e pela disponibilização dos dados de processo de forma modular e interoperável.
    \item Implementar a interface web de monitoramento, permitindo a visualização dos principais indicadores, tendências ou telas resumo associadas ao processo acompanhado.
    \item Validar o módulo de monitoramento por meio de testes e de um cenário demonstrativo (ou estudo de caso), evidenciando o correto acesso às variáveis e a estabilidade da solução frente ao volume de dados esperado.
\end{itemize}